\documentclass[11pt]{article}

\usepackage[margin=1in]{geometry}
\usepackage[T1]{fontenc}
\usepackage{lmodern}
\usepackage{fontspec}
\usepackage{amsmath}
\usepackage{amssymb}
\usepackage{booktabs}
\usepackage{longtable}
\usepackage{array}
\usepackage{ragged2e}
\usepackage{xcolor}
\PassOptionsToPackage{hyphens}{url}
\usepackage[colorlinks=true,linkcolor=blue!50!black,citecolor=blue!50!black,urlcolor=blue!50!black]{hyperref}
\usepackage[round]{natbib}
\usepackage{microtype}
\usepackage{enumitem}

% This paper's own section numbers (e.g. "4.2", "4.3", "4.4" -- Section 4
% deliberately has no "4.1"; Section 8's "8.3.1" deliberately appears before
% "8.1") are hard-coded as literal text in every heading. Disabling LaTeX's
% automatic numbering here (rather than "fixing" the numbering) is the
% format-only-preserving choice -- renumbering would be a content change,
% not a formatting one.
\setcounter{secnumdepth}{-5}

% Pandoc's tight-list marker (used throughout the converted body).
\providecommand{\tightlist}{%
  \setlength{\itemsep}{0pt}\setlength{\parskip}{0pt}}

\title{Tool-Description Quality Is Not One Axis: A Regime Analysis of Where It Helps and Where It Backfires}
\author{}
\date{}

\begin{document}

\maketitle

\begin{abstract}
Tool-description quality is widely treated as a broadly-applicable lever
for agent tool-use, but it is not a single better/worse axis: the
precision that helps an agent disambiguate within a family of confusable
tools is orthogonal to, or actively harmful for, context-rich selection
and for tool retrieval. We test this with a single frozen evaluation
protocol --- one classifier, one judge, one generator family,
pre-registered thresholds --- across a synthetic confusable-catalog
experiment, two real production MCP-server mirrors (GitHub, AWS IAM), a
synthetic internal-proxy catalog, and a pre-registered pilot of ten
public Python MCP servers. The effect is real but regime-bounded, not a
general law: an oracle description helps at one tested catalog density
(60 tools/10 families, +34.5pp on gemma2:9b, and not collapsing on a
substantially stronger model, Llama-3.3-70B) with no headroom; realizing
it safely through automatic generation is a separate condition,
requiring documented source. Outside these conditions the effect is
null or reverses (zero headroom on two well-documented production
servers; −20pp harm on one already-resolved family; harm to retrieval
across three retriever types). Contributions: a falsifiable regime map
of where description quality helps, harms, or does nothing; a
pre-registered prevalence measurement finding the behavioral regime in
0 of 9 testable Python MCP servers --- a lower bound, not a population
estimate; and a localizability boundary --- a pairwise LLM-judge
confusability method fails under two independent framings, via two
distinct failure mechanisms.

\end{abstract}

\hypertarget{introduction}{%
\section{1. Introduction}\label{introduction}}

This paper's central finding is counterintuitive enough to state up
front, with its scope attached: on the fixtures tested here,
tool-description quality is not a single better/worse axis. The same
behavioral-axis precision that helps an agent disambiguate within a
family of confusable tools (§4.2.1, one density point: 60 tools / 10
families) is not what a context-rich agent already resolving from task
wording needs, and is actively harmful to tool retrieval (§4.3.3, one
synthetic catalog, three retriever types). Optimizing a description for
one of these objectives can measurably neutralize or harm another --- a
real, mechanistically legible pattern on the fixtures tested, not a
demonstrated general law about description quality. The rest of this
paper places that finding inside a broader map (Section 4), asks how
common the map's ``it matters'' region is in real servers (Section 5),
and tests whether that region can be found cheaply before it is needed
(Section 6).

\hypertarget{the-assumption-under-test}{%
\subsection{1.1 The assumption under
test}\label{the-assumption-under-test}}

Agentic systems that call external tools --- via MCP servers,
OpenAI-style function schemas, or similar interfaces --- depend on an
LLM agent correctly selecting which tool to invoke for a given task.
Practitioner-facing engineering guidance treats the \emph{quality of the
tool description text} as a primary lever for improving that selection:
Anthropic's own tool-use engineering guidance states that ``even small
refinements to tool descriptions can yield dramatic improvements,''
reporting that precise refinements to tool descriptions helped Claude
Sonnet 3.5 reach state-of-the-art on SWE-bench Verified (Anthropic
Engineering, ``Writing tools for agents'' --- publication date not
independently confirmed during this paper's preparation, cited for content
only) \nocite{anthropic2026tools}. GitHub's
engineering team, running an offline evaluation harness against their
own MCP server --- the same \url{github/github-mcp-server} mirrored
in RW1 (§4.3.1) --- reports the same sensitivity from the other side:
``tightening a description, adding or removing a tool, or combining a
few similar tools can shift results a lot'' (GitHub Engineering,
``Measuring what matters: how offline evaluation of GitHub MCP Server
works'' --- publication date not independently confirmed during this paper's preparation,
cited for content only) \nocite{github2026measuring}. This appears to sit in tension with RW1's own
finding of zero headroom on that exact server for gemma2:9b (§4.3.1) ---
we do not resolve this tension here, since GitHub's harness details
(which model(s), which metric, which task distribution) were not
verified during this paper's preparation beyond the quoted sentence; flagged as an open
question rather than silently juxtaposed as if the two results agree. An
empirical audit of 856 tools across 103 real MCP servers found 97.1\% of
tool descriptions contain at least one quality ``smell'' and 56\% fail
to state their purpose clearly, and that LLM-based description
augmentation lifts task success by a median of +5.85pp --- but also
regresses performance in 16.67\% of cases and inflates execution steps
by 67.46\% (Hasan et al., ``Model Context Protocol (MCP) Tool
Descriptions Are Smelly!,'' arXiv:2602.14878, 2026) \nocite{hasan2026smelly}. Real tool-name
collisions cause user-visible failures in production agent frameworks,
independent of any synthetic fixture (e.g., a reported GitHub issue:
``Duplicate tool names across MCP servers cause errors,''
\url{openai/openai-agents-python#464}) \nocite{openai2024issue464}. Taken together, this is real
evidence that description quality matters \emph{somewhere} and that even
small edits move outcomes --- but none of these sources characterize the
\emph{boundary} of where it matters, whether it generalizes across
catalog scales, or whether the same fix that helps selection also helps
or harms retrieval and already-correctly-resolving cases. That is the
gap this paper addresses.

\hypertarget{what-this-paper-does-instead}{%
\subsection{1.2 What this paper does
instead}\label{what-this-paper-does-instead}}

We hold three properties fixed across every experiment reported here:
one frozen evaluation protocol (Section 3), one classifier, and
null-first-class reporting --- an aborted or negative result is reported
as a finding, not omitted. Against that fixed protocol we ask three
questions in sequence:

\begin{enumerate}
\def\labelenumi{\arabic{enumi}.}
\tightlist
\item
  \textbf{Where does description quality change agent behavior, and
  where doesn't it?} (Section 4, EXP-4 --- a regime map assembled from
  every experiment run under the frozen protocol.)
\item
  \textbf{How common is the ``it matters'' regime in real, sampled MCP
  servers?} (Section 5, EXP-1 --- a behavioral prevalence measurement on
  a pre-registered sample.)
\item
  \textbf{Can the regime be found cheaply, before deploying an agent
  against a real server?} (Section 6, EXP-3 --- a pairwise confusability
  localizer, tested against behavioral ground truth.)
\end{enumerate}

\hypertarget{contributions}{%
\subsection{1.3 Contributions}\label{contributions}}

\begin{itemize}
\tightlist
\item
  \textbf{A falsifiable, regime-bounded map} of where tool-description
  quality changes agent selection behavior. Two distinct conditions
  compose it, from different experiments: (i) a hand-written oracle
  \emph{effect}, bracketed but not precisely located on the density axis
  (Section 4.2.1);

  \begin{enumerate}
  \def\labelenumi{(\roman{enumi})}
  \setcounter{enumi}{1}
  \tightlist
  \item
    realizing that effect safely through \emph{automatic generation}
    additionally requires documented source code (Section 4.2.2).
    Outside this intersection the effect is null or negative (Section
    4.3).
  \end{enumerate}
\item
  \textbf{A prevalence measurement}: on a pre-registered pilot sample of
  10 Python-only public MCP servers, 0 of 9 servers with a testable
  confusable family showed in-regime behavior, using the general
  behavioral regime definition (§5.1) --- a lower bound, not a closed
  claim, on how rare the regime is among public, documented servers
  (Section 5).
\item
  \textbf{A localizability boundary}: a pairwise LLM-judge confusability
  method --- the natural cheap alternative to running the full
  behavioral protocol against every server --- fails under two
  independently pre-registered framings, reaching the identical
  confusion matrix both times via two different degeneracy mechanisms
  rather than one repeated failure mode (Section 6). This extends the
  map from \emph{where the regime occurs} to \emph{where it can, and
  cannot, be found cheaply}.
\end{itemize}

\hypertarget{roadmap}{%
\subsection{1.4 Roadmap}\label{roadmap}}

Section 2 positions this work against tool-use benchmarks, MCP
documentation practice, and retrieval-augmented tool selection. Section
3 fixes the evaluation protocol every experiment reuses. Section 4
presents the regime map: where the effect helps, and --- as part of the
same map, not a separate cautionary appendix --- where it does nothing
or actively harms. Section 5 presents the prevalence result. Section 6
presents the localizer result. Section 7 discusses the synthesis.
Sections 8--9 cover threats to validity and the reproducibility
artifact.

\begin{center}\rule{0.5\linewidth}{0.5pt}\end{center}

\hypertarget{related-work}{%
\section{2. Related Work}\label{related-work}}

\hypertarget{tool-use-and-function-calling-benchmarks}{%
\subsection{2.1 Tool-use and function-calling
benchmarks}\label{tool-use-and-function-calling-benchmarks}}

Benchmarks for LLM tool/function-calling typically measure whether an
agent calls the \emph{correct} tool and constructs a \emph{valid} call,
usually against a fixed, hand-built catalog. Selection reliability is
known to degrade as the tool menu grows --- a synthetic causal-filtering
study reports wrong-tool calls increasing with larger tool menus across
four LLM backends (Babu and Iyer, ``ToolChoiceConfusion: Causal Minimal
Tool Filtering for Reliable LLM Agents,'' arXiv:2606.06284,
2026) \nocite{babu2026toolchoice}, and a
large real-usage benchmark of tool-use ``in the wild'' finds no model,
of 57 tested, exceeds 15\% accuracy on tasks reflecting authentic
multi-turn, mixed-intent user behavior (Yu et al., ``Benchmarking LLM
Tool-Use in the Wild'' {[}WildToolBench{]}, arXiv:2604.06185,
2026) \nocite{yu2026wildtool}.
Neither benchmark isolates description \emph{quality} specifically as
the independent variable --- which is why the regime map in Section 4
spans purpose-built synthetic catalogs at controlled scale (T18)
alongside real production-server mirrors (RW1, RW2) rather than a single
fixture designed for a different question.

\hypertarget{mcp-ecosystem-and-documentation-practice}{%
\subsection{2.2 MCP ecosystem and documentation
practice}\label{mcp-ecosystem-and-documentation-practice}}

The Model Context Protocol ecosystem has produced a growing set of
public servers with widely varying documentation quality (Hasan et al.,
2026, above \nocite{hasan2026smelly}: 97.1\% of a 856-tool sample shows at least one description
``smell''), and a parallel set of tools (linters, description scorers,
\texttt{llms.txt} conventions) that assume documentation quality is both
measurable and improvable in a way that changes agent behavior. Section
4.3.4 and Section 6 directly test the measurability half of that
assumption for one specific construct --- can a scorer identify
\emph{which tool pairs} in a catalog are confusable --- and find it
fails at both the single-score and pairwise level as currently
constructed, on the fixtures tested here.

\hypertarget{retrieval-augmented-tool-selection}{%
\subsection{2.3 Retrieval-augmented tool
selection}\label{retrieval-augmented-tool-selection}}

Where a catalog is too large to place in-context, tool retrieval
(lexical or embedding-based) is used to narrow the candidate set before
selection. A large retrieval benchmark (7.6k tasks, 43k tools) finds
even strong information-retrieval models perform poorly at this task ---
the best-performing model reaches only 33.83 nDCG@10 (Shi et al.,
``Retrieval Models Aren't Tool-Savvy: Benchmarking Tool Retrieval for
Large Language Models'' {[}ToolRet{]}, ACL 2025 Findings /
arXiv:2503.01763) \nocite{shi2025toolret}. One proposed remedy is LLM-based description
\emph{expansion}, reported to yield state-of-the-art retrieval gains on
two benchmarks (Lu et al., ``Tools are under-documented: Simple Document
Expansion Boosts Tool Retrieval'' {[}Tool-DE{]}, arXiv:2510.22670,
2025) \nocite{lu2025toolde}. Retrieval and direct-selection description quality are often
assumed to move together on this basis --- richer descriptions should
help both. Section 4.3.3 shows they are anti-correlated \emph{on the one
synthetic catalog tested here}: the same behavioral-axis precision that
helps direct-selection disambiguation matches poorly against the coarse,
intent-level queries a retrieval index is queried with --- a result that
sits in tension with, rather than replicating, Tool-DE's
description-expansion gains, and that we do not claim generalizes beyond
this fixture (see §8.1).

\hypertarget{positioning}{%
\subsection{2.4 Positioning}\label{positioning}}

This paper positions against the implicit claim, common across the
practices surveyed above, that tool-description-quality improvement is a
broadly-applicable, low-risk intervention. We do not dispute that it
helps in \emph{some} regime --- Section 4.2 characterizes one such
regime, bounded by the specific fixtures tested --- but we show that
regime is narrower and less precisely located than a flat ``improve your
descriptions'' recommendation implies, and that the same intervention is
risky outside it.

\begin{center}\rule{0.5\linewidth}{0.5pt}\end{center}

\hypertarget{methods-the-frozen-evaluation-protocol}{%
\section{3. Methods --- The Frozen Evaluation
Protocol}\label{methods-the-frozen-evaluation-protocol}}

All results in this paper are produced under a single protocol,
committed once (\url{docs/research/frozen_protocol.md}) and never
edited after any experiment's results were collected. This section
states it in full; individual experiment sections do not repeat it.

\hypertarget{governance}{%
\subsection{3.1 Governance}\label{governance}}

Every experiment's spec --- task set, gold labels, stability-screen
procedure, and sidedness of the sign test --- is pre-registered to its
branch \textbf{before} any run starts. Any change to the judge, scorer,
rubric, or calibration constants (``condition \#1'') is escalated as a
draft PR for human review before merge. Null and aborted results are
reported as-is; results are never tuned after being observed.
Judge, generator, and agent model families are structurally independent
of any assistant used to author or orchestrate this research
program --- no shared credentials, infrastructure, or model family.

\hypertarget{frozen-configuration}{%
\subsection{3.2 Frozen configuration}\label{frozen-configuration}}

\renewcommand{\arraystretch}{1.15}
\begin{longtable}{@{}p{0.26\textwidth}p{0.68\textwidth}@{}}
\toprule
\textbf{Component} & \textbf{Value} \\
\midrule
\endhead
Judge model / seed & \texttt{llama3.1:8b} / \texttt{42} \\
Generator model & \texttt{qwen3:8b} (always distinct from the judge) \\
Default agent & \texttt{gemma2:9b} \\
Trials per arm & 5 (default); \textbf{3} for FRONTIER-T18 and EXP-3's
judge calls --- each independently pre-registered and justified (API
rate/cost limits for FRONTIER-T18; matching the codebase's existing
graded-judge convention for EXP-3) \\
Sign-test $\alpha$ & 0.05, two-sided unless pre-registered one-sided \\
Headroom ceiling & 0.85 --- if Arm A (unmodified descriptions) already
scores $\geq$85\% parse-success accuracy on the contested set, the experiment
aborts as a null (no headroom), rather than proceeding to an A/B
comparison \\
Minimum contested tasks & 6 \\
\bottomrule
\end{longtable}
\renewcommand{\arraystretch}{1}

\hypertarget{classifier}{%
\subsection{3.3 Classifier}\label{classifier}}

Every trial is classified into exactly one of three outcomes, plus a
separate flag:

\begin{itemize}
\tightlist
\item
  \texttt{SELECTED-CORRECT} --- the agent selected the pre-registered
  gold tool.
\item
  \texttt{SELECTED-WRONG} --- the agent selected a different tool.
\item
  \texttt{ABSTAINED-OR-HEDGED} --- the agent produced output but hedged
  or explicitly abstained.
\item
  \texttt{PARSE-FAILED} --- the output could not be parsed into a
  structured selection at all.
\end{itemize}

\texttt{PARSE-FAILED} is always reported and never folded into the
three-outcome denominator; a model or condition that appears to show
``no effect'' is checked against its parse-failure rate before that
conclusion is accepted (Section 8.3).

\hypertarget{effect-measurement}{%
\subsection{3.4 Effect measurement}\label{effect-measurement}}

The unit of analysis is the \textbf{task}, not the trial. Arm A is run
twice with independent seeds as a stability screen; a task whose per-run
accuracy differs by more than one trial is dropped from the primary
analysis and reported separately. Effect is Arm B accuracy minus Arm A
accuracy on parse-success, stable, contested tasks; the sign test is
computed per-task (B-beats-A vs.~A-beats-B vs.~tie), not per-trial.

\hypertarget{cross-experiment-comparability}{%
\subsection{3.5 Cross-experiment
comparability}\label{cross-experiment-comparability}}

Cross-experiment effect-size comparisons (e.g., ``does the T18 effect
hold at a different capability tier?'') are valid only when the fixture,
judge, classifier, and sign-test procedure are identical and the
\emph{only} varying factor is the one under test. The T18 gemma2:9b
(+34.5pp) and FRONTIER-T18 Llama-3.3-70B (+40.8pp) figures reported in
Section 4.2.3 are explicitly \textbf{not} part of a single controlled
ladder --- they come from different harnesses and hosts --- and are
never combined into a claim that the effect ``grows with capability,''
only that it ``survives'' at a higher capability tier tested
independently. EXP-2, the controlled same-family capability ladder that
would test this properly, was dropped from this paper's scope (see
Section 8.2).

\begin{center}\rule{0.5\linewidth}{0.5pt}\end{center}

\hypertarget{exp-4-regime-map}{%
\section{4. EXP-4 --- Regime Map}\label{exp-4-regime-map}}

\hypertarget{what-exp-4-consolidates}{%
\subsection{4.1 What EXP-4 consolidates}\label{what-exp-4-consolidates}}

EXP-4 makes no new runs; it consolidates every experiment already run
under the frozen protocol into a single map of where tool-description
quality changes agent behavior and where it does not. Framing A treats
this map, together with the prevalence finding in Section 5, as the
paper's primary artifact --- the field's assumption is not ``false,'' it
is \textbf{regime-bounded}, and this section states the boundary
precisely enough to be checked against a new server.

\hypertarget{where-it-helps}{%
\subsection{4.2 Where it helps}\label{where-it-helps}}

\hypertarget{confusable-at-scale-t18}{%
\subsubsection{4.2.1 Confusable-at-scale
(T18)}\label{confusable-at-scale-t18}}

On a synthetic 60-tool catalog (10 families of 6 near-neighbor tools
each, gemma2:9b agent), empty descriptions leave the agent to resolve
tool selection from names alone: parse-success accuracy on the contested
subset is \textbf{62.9\%} (110/175). Oracle descriptions that target
each family's within-family distinguishing dimension (storage medium,
operation scope, permanence, notification channel, computation type)
raise this to \textbf{97.4\%} (190/195) --- a \textbf{+34.5pp}
discrimination effect (task-clustered sign test: n+=16, n−=0, ties=24,
p\textless0.0001). A second, mechanistically distinct effect appears
alongside it: at this catalog density, empty descriptions also
destabilize call \emph{formation} itself --- parse-failure rate falls
from 12.5\% (25/200) to 2.5\% (5/200) under oracle descriptions,
contributing roughly 5--6pp of the aggregate delta separately from the
discrimination effect.

This effect appears \textbf{density-gated} rather than a general
property of description quality, but the paper has only two points on
the density axis, not a measured curve: an earlier fixture at 16 tools /
8 clusters (T17) placed Arm A at 81.2\% accuracy from names alone ---
headroom too high (by the ad hoc 70\% target that experiment used) for
the oracle arm to be run at all, so the description effect was never
actually tested at that density, not tested-and-absent. The
60-tool/10-family point is positive (+34.5pp); the 16-tool/8-cluster
point is untested; nothing between 16 and 60 tools has been run. The
defensible claim is that density gates \emph{whether the effect is even
testable} (headroom exists or doesn't) at these two specific fixtures
--- not that density is \emph{the} gating variable in general, and not
that the transition point is located more precisely than ``somewhere in
this untested bracket.''

\hypertarget{under-documented-source-with-docstrings-q3-q6-guard-b}{%
\subsubsection{4.2.2 Under-documented source with docstrings (Q3 → Q6,
Guard-B)}\label{under-documented-source-with-docstrings-q3-q6-guard-b}}

The T18 result establishes that oracle (hand-written, ground-truth)
descriptions help; it does not establish that a \emph{generator} can
produce them safely. A four-stage progression (Q3, Q4, Q5, Q6) on a
12--23-tool fixture establishes both the recovery and safety conditions;
the two endpoints carry the argument (full intermediate progression ---
including two distinct fabrication failure modes surfaced and closed
along the way --- in Appendix A.7):

\renewcommand{\arraystretch}{1.15}
\begin{longtable}{@{}p{0.34\textwidth}p{0.20\textwidth}p{0.12\textwidth}p{0.24\textwidth}@{}}
\toprule
\textbf{Condition} & \textbf{Recovery (6 structural contested tasks)} & \textbf{p} &
\textbf{No-fabrication} \\
\midrule
\endhead
Interface-only, no source (Q2a/Q2b) & 12.5\% (11.1pp of 88.9pp
available) & 0.50 (n.s.) & n/a --- nothing to fabricate from \\
Q5/Q6-guarded: scoped source + docstrings + target-grounded prompt (Q6 =
23-tool blanket application) & 100\% (6/6) & 0.0313 (Q5) / 0.0312 (Q6) &
PASS, \textbf{0/11 regressions} on already-passing tools (Q6) \\
\bottomrule
\end{longtable}
\renewcommand{\arraystretch}{1}

Two findings compose into the deployment answer. First,
\textbf{generation from the interface alone cannot recover the T18 gap}:
two independent interface-only generation strategies (per-tool, Q2a;
catalog-aware, Q2b) both recover only 12.5\% of the oracle gain (11.1pp
of 88.9pp available), not significant (p=0.50). The T18-decisive
distinctions live in tool \emph{behavior} --- absent from both tool
names and identical parameter schemas --- and no amount of
interface-level prompting recovers information that is not present in
the interface. Second, \textbf{source access alone is not automatically
safe}: whole-file source with docstrings stripped, and per-tool scoped
source with docstrings present, each open a distinct fabrication failure
mode (cross-tool symbol misattribution; docstring-body vocabulary gaps,
respectively). The safe \emph{and} fully recovering configuration ---
scoped source, docstrings retained, target-grounded prompting that
forbids comparative neighbor claims (Guard-B) --- was reached only at
the fourth iteration (Q5), and confirmed to cause zero regressions when
applied blanket across a 23-tool catalog including three collision-prone
pairs (Q6).

\hypertarget{strong-agent-survival-frontier-t18}{%
\subsubsection{4.2.3 Strong-agent survival
(FRONTIER-T18)}\label{strong-agent-survival-frontier-t18}}

A natural objection to Section 4.2.1 is that the effect might be an
artifact of gemma2:9b's specific capability level --- a stronger agent
might resolve the same catalog from names alone. We re-ran the identical
T18 fixture and oracle descriptions against Llama-3.3-70B (OpenRouter),
a substantially stronger open-weight agent, using a 3-outcome classifier
and 3 trials per task. Arm A accuracy was \textbf{59.2\%} (71/120) ---
comparable headroom to the gemma2:9b run --- and Arm B (oracle) reached
\textbf{100.0\%} (120/120): an effect of \textbf{+40.8pp} (sign test
n+=19, n−=0, ties=21, p\textless0.0001; on the stable set excluding 5
Arm-A trial-flippers, n+=14, n−=0, p\textless0.001). All 19 Arm-A miss
tasks were fully recovered by the oracle, with no oracle-resistant floor
in this fixture. The effect does not collapse, and does not shrink,
moving from a 9B to a 70B agent.

This is a survival claim, not a growth claim, and not a frontier claim:
Llama-3.3-70B is one strong open model, one data point, not a
Claude/GPT-class proprietary frontier model, and the +40.8pp figure is
\textbf{not directly comparable} to T18's +34.5pp --- different harness,
different classifier host (Section 3.5). The result is
\textbf{inconsistent with} ``the T18 effect is a gemma2:9b-specific
capability artifact'' as an explanation; it is one additional data
point, on a different harness, not a proof that no capability tier would
show the effect shrink --- and it does not close the question of whether
a true proprietary frontier model would behave differently.

\hypertarget{where-it-doesnt-help-or-actively-harms}{%
\subsection{4.3 Where it doesn't help, or actively
harms}\label{where-it-doesnt-help-or-actively-harms}}

This is not a separate cautionary appendix --- under Framing A it is the
other half of the same map, and the boundary in Section 4.3 is only
precise because these non-regimes were tested under the identical
protocol.

\hypertarget{no-headroom-real-servers-rw1-rw2}{%
\subsubsection{4.3.1 No-headroom real servers (RW1,
RW2)}\label{no-headroom-real-servers-rw1-rw2}}

Two production MCP-server mirrors --- GitHub
(\url{github/github-mcp-server}, 21 tools, 5 confusable families,
real docstrings) and AWS IAM (\url{awslabs/mcp}, 29 tools, 3
confusable families, including four-way attach/detach/user/group
confusability) --- were tested with real, unmodified docstrings. Both
resolved \textbf{100\%} of tasks correctly, including every
pre-registered contested task (21/21 for GitHub; 29/29 including 12
contested for AWS IAM). There is no headroom for a description fixer to
recover on either server: gemma2:9b resolves these tools from
task-context wording alone, and the terse, one-sentence docstrings both
servers already ship are sufficient. The regime identified in Section
4.2.1--4.2.2 requires headroom that these two real, well-documented
servers simply do not have.

\hypertarget{harm-on-an-already-resolved-family-p2-a-account_query}{%
\subsubsection{4.3.2 Harm on an already-resolved family (P2-A
account\_query)}\label{harm-on-an-already-resolved-family-p2-a-account_query}}

On a 48-tool synthetic internal-proxy catalog built to test whether
description improvement is uniformly safe across many confusable
families, one family (\texttt{account\_query}, 5 tools) shows the
opposite of the T18/Q-series pattern. Under thin descriptions, the
task's SQL-like WHERE-clause phrasing heuristically matches
\texttt{query\_accounts} by name: \textbf{100\%} (5/5) correct. Under
both Guard-B-generated and hand-written oracle descriptions, more
precise phrasing (``Executes a parameterized SQL-like WHERE clause'')
introduces disambiguation noise among five semantically related tools,
and accuracy \textbf{drops to 80\%} (4/5) under both --- a \textbf{−20pp
harm}, reproducing at the same magnitude under two independent
generation strategies, ruling out noise as the explanation. Two other
families on the same catalog (\texttt{ticket\_lifecycle}:
delete/purge/archive/ expire; \texttt{invoice\_write}:
update/upsert/patch/replace/amend) stay at 100\% across all three arms
--- task-context phrasing already signals the dangerous distinction
(``this cannot be undone''), and better descriptions add no incremental
safety there either way. The one family that is \emph{fully} recovered
from a genuine failure (\texttt{order\_read}, 0\%→100\% under both
Guard-B and oracle) is a low-stakes return-shape disambiguation, not a
high-stakes one. Aggregate across all 31 contested tasks (A=77.4\%,
Guard-B=96.8\%, Oracle=93.5\%) is directional but underpowered (sign
test p=0.07/0.125); the account\_query harm is the reproducible,
per-family finding, and the aggregate number should not be read as a
headline recovery rate.

\hypertarget{retrieval-readiness-is-anti-correlated-not-neutral-f2}{%
\subsubsection{4.3.3 Retrieval-readiness is anti-correlated, not neutral
(F2)}\label{retrieval-readiness-is-anti-correlated-not-neutral-f2}}

If a catalog is too large to place in-context, tool selection is often
mediated by a retrieval step (lexical or embedding-based lookup against
description text) before the agent ever sees a candidate set. Testing
the same P2-A catalog's descriptions as retrieval documents against 93
underspecified natural-language queries, across three retriever types:

\renewcommand{\arraystretch}{1.15}
\begin{table}[h]
\centering
\begin{tabular}{@{}lccc@{}}
\toprule
\textbf{Arm} & \textbf{BM25 MRR} & \textbf{TF-IDF MRR} & \textbf{Embedding MRR (nomic-embed-text)} \\
\midrule
Thin & \textbf{0.304} & \textbf{0.317} & \textbf{0.510} \\
Guard-B & 0.225 ($-$0.079) & 0.204 ($-$0.113) & 0.286 ($-$0.224) \\
Oracle & 0.234 & 0.228 & 0.362 \\
\bottomrule
\end{tabular}
\end{table}
\renewcommand{\arraystretch}{1}

On this catalog, thin descriptions \textbf{outperform} Guard-B and
oracle descriptions on every retriever tested, with the harm larger
under semantic embedding retrieval than lexical retrieval. The mechanism
is legible: compact verb-noun descriptions (``Get the order.'') keyword-
and semantically-match coarse, intent-level queries (``get an order'')
directly; precision-optimized descriptions (``Returns order summary
fields including status, total, item count\ldots{}'') match at the
implementation level --- exactly the level that helps within-family
selection disambiguation (Section 4.2.1) and hurts coarse retrieval
matching here. The property that makes a description good for one
downstream use (direct selection in a fully-listed catalog) makes it
worse for another (retrieval against a partial query) on this synthetic
proxy --- this is the central non-obvious finding connecting Sections
4.2 and 4.3, tested on one fixture and not yet checked against a real
server's retrieval behavior (§8.1).

\hypertarget{single-score-judging-cannot-localize-which-pairs-are-confusable}{%
\subsubsection{4.3.4 Single-score judging cannot localize which pairs
are
confusable}\label{single-score-judging-cannot-localize-which-pairs-are-confusable}}

Both real-server mirrors were also scored with AgentGauge's existing
\texttt{discoverability} judge, whose sub-score returns a single number
per catalog. On GitHub, the score is flat at \textasciitilde70/100
across every family, and does not flag either of the two families
GitHub's own maintainers consolidated to reduce confusion (0/2 overlap).
On AWS IAM, the score is flat at \textasciitilde68.7/100, and does not
flag any of the three contested families (0/3 overlap); the accompanying
name-collision heuristic instead flags verb-antonym pairs
(\url{attach_user_policy}/\url{detach_user_policy}) that
caused zero real confusion. A single number per catalog is structurally
incapable of naming \emph{which} pair is the problem --- this is the
motivating gap for the localizer tested in Section 6, not a calibration
issue that more judge tuning would fix.

\hypertarget{description-quality-is-multi-objective}{%
\subsection{4.4 Description quality is
multi-objective}\label{description-quality-is-multi-objective}}

This is the paper's central synthesis, previewed in the Introduction and
Abstract: the same behavioral-axis precision that helps within-family
selection is not one general ``better'' direction for description
quality --- on the fixtures tested here, it is orthogonal or actively
harmful to other things a description interacts with (what a
context-rich agent needs; what retrieval rewards). Every hedge below is
load-bearing, not decorative: each line names exactly which fixture the
claim is scoped to.

\begin{verbatim}
Description precision HELPS within-family discrimination when (each tested at
one fixture/point, not a swept curve):
 - Catalog density is at least as high as the one positive point tested
   (§4.2.1), AND
 - Source access includes docstrings (safe generation requires this;
   body-only source opens a fabrication vector), AND
 - Arm A shows real headroom (the agent is not already resolving from
   task-context wording).

The same precision is ORTHOGONAL or ANTI-CORRELATED to, on the fixtures tested:
 - what context-rich agents need (RW1, RW2, and P2-A's high-stakes
   families all resolve at 100% from context alone, regardless of
   description quality),
 - what retrieval rewards on one synthetic catalog (F2: harm across all
   three retriever types tested there), and
 - what a single-score judge can report (localization requires a
   pairwise comparison, tested next in Section 6).
\end{verbatim}

\begin{center}\rule{0.5\linewidth}{0.5pt}\end{center}

\hypertarget{exp-1-server-population-prevalence}{%
\section{5. EXP-1 --- Server-Population
Prevalence}\label{exp-1-server-population-prevalence}}

\hypertarget{question-and-regime-definition}{%
\subsection{5.1 Question and regime
definition}\label{question-and-regime-definition}}

Section 4 establishes \emph{where} the effect exists; it does not
establish \emph{how common} that regime is among real MCP servers. We
define ``in-regime'' behaviorally, not by any static description-quality
property (which would be circular): a confusable family on a server is
in-regime iff, under the frozen protocol, (a) the server's real shipped
descriptions fail at least one contested task, \textbf{and} (b) an
oracle description recovers it. ``Thin descriptions'' is an input
property; only behavior the fix actually repairs counts as in-regime.

\textbf{This is the general behavioral construct, not a re-test of
Section 4.2.1's specific density point.} None of the ten sampled servers
were selected or verified to actually reach the 60-tool/10-family
density where T18 found an effect (most have far fewer tools; the one
large-catalog server sampled, at 116 tools, failed for an unrelated
reason --- catalog overwhelm, malformed output under a large listing,
not confident wrong-tool selection --- and contributes no evidence
either way about the density regime specifically). What EXP-1 measures
is the prevalence of the general two-condition behavioral regime across
whatever confusable families each server happens to have, at whatever
scale it happens to be built at --- a broader and different question
than ``how many real servers are built at T18's specific tested
density.''

\hypertarget{sampling-frame}{%
\subsection{5.2 Sampling frame}\label{sampling-frame}}

The pre-registered, ratified frame is \textbf{N=10, Python-only,
doc-density-stratified}, sourced from public GitHub; 11 non-Python
servers were dropped because the generic regex fallback extractor used
for non-Python source proved unreliable on them (full revision history
in Appendix A.6). This is a narrower claim than ``public MCP servers''
--- it is ``Python MCP servers on GitHub, N=10 pilot'' --- stated as
such everywhere this number appears in this paper.

\hypertarget{result}{%
\subsection{5.3 Result}\label{result}}

\begingroup
\footnotesize
\renewcommand{\arraystretch}{1.2}
\begin{longtable}{@{}p{0.20\textwidth}p{0.10\textwidth}p{0.13\textwidth}p{0.13\textwidth}p{0.07\textwidth}p{0.27\textwidth}@{}}
\toprule
\textbf{Server} & \textbf{Tier} & \textbf{Arm A} & \textbf{Arm B (oracle)} & \textbf{Effect} & \textbf{Verdict} \\
\midrule
\endhead
github-mcp (RW1, anchor) & anchor & 100\% (21/21) & --- & n/a &
OUT-OF-REGIME \\
aws-iam-mcp (RW2, anchor) & anchor & 100\% (29/29 incl.\ 12 contested) &
--- & n/a & OUT-OF-REGIME \\
lucasastorian-llmwiki & near\_empty & 100\% & --- & n/a &
OUT-OF-REGIME \\
stefanoamorelli-sec-edgar-mcp & thin & 100\% & --- & n/a &
OUT-OF-REGIME \\
stickerdaniel-linkedin-mcp-server & thin & 100\% & --- & n/a &
OUT-OF-REGIME \\
mrexodia-ida-pro-mcp & near\_empty & 90\% & --- & n/a & OUT-OF-REGIME
(above headroom ceiling) \\
AminForou-mcp-gsc & well\_documented & 75\% & 75\% & +0pp &
OUT-OF-REGIME (real functional overlap; no description fixes it) \\
taylorwilsdon-google\_workspace\_mcp & thin & 0\% & 0\% & +0pp &
OUT-OF-REGIME (catalog-overwhelm failure mode, not confusability) \\
datalayer-jupyter-mcp-server & near\_empty & 70\% & 55\% &
\textbf{$-$15pp} & OUT-OF-REGIME (HARM --- replicates the P2-A
account\_query pattern on a fresh real server) \\
Dataojitori / blazickjp-arxiv / LycheeMem & well\_documented & --- & ---
& --- & no testable confusable family (3 servers) \\
\bottomrule
\end{longtable}
\endgroup

\textbf{Headline: 0 of 9 servers with a testable confusable family show
in-regime behavior} (7 freshly scored + 2 RW1/RW2 anchors cited from
Section 4.3.1). Per fresh-only tier: well\_documented 0/1 testable, thin
0/3, near\_empty 0/3 --- a clean null across every tier of this pilot.
Three of the ten servers had no genuinely confusable family at all
(mechanical clustering found none, or the one candidate was rejected on
manual review as not genuinely confusable).

\hypertarget{the-seed-bug-episode}{%
\subsection{5.4 The seed-bug episode}\label{the-seed-bug-episode}}

A first trial run passed a single fixed seed (\texttt{seed=42}) to every
one of five nominal repeated trials per task, instead of the codebase's
established \texttt{seed=42+trial\_idx} convention --- meaning the first
run sampled zero real trial-to-trial variance. That run reported two
servers as in-regime (\texttt{mrexodia-ida-pro-mcp} +50pp;
\texttt{datalayer-jupyter-mcp-server} +25pp). Fixing the seed and
re-running \textbf{reversed both findings}: ida-pro's real Arm A
accuracy is 90\% (correctly aborts, no headroom), and jupyter's Arm B is
a genuine harm (−15pp, not a +25pp recovery). Both apparent in-regime
results were seed artifacts, caught before being reported as findings
--- a credibility asset for the pipeline's own error-detection, reported
here as such.

This asset has a mirror-image limitation that is \emph{more} important
than the credit above, and is given its own standalone treatment, not
folded into this paragraph: \textbf{the note opening §8} states why a bug that
suppressed a real in-regime signal, rather than manufacturing a false
one, would not have been caught the same way.

\hypertarget{scope}{%
\subsection{5.5 Scope}\label{scope}}

Public servers skew documented relative to the internal/custom long
tail; this prevalence figure is a \textbf{lower bound} on how common the
regime is, not a closed population estimate. We do not claim the regime
``does not occur'' --- we claim it is rare in this sampled population of
public, source-available Python MCP servers as of the frame's fixed
date, and that the under-documented internal segment this bounds is, by
construction, unsampleable from public data.

\begin{center}\rule{0.5\linewidth}{0.5pt}\end{center}

\hypertarget{exp-3-confusability-localization}{%
\section{6. EXP-3 --- Confusability
Localization}\label{exp-3-confusability-localization}}

\hypertarget{motivation}{%
\subsection{6.1 Motivation}\label{motivation}}

Section 4.3.4 established that the existing single-score
\texttt{discoverability} judge cannot say \emph{which} tool pair in a
catalog is confusable --- a structural limit of returning one number per
catalog. If the regime Section 4 characterizes is to be found cheaply,
before running the full behavioral protocol against every server, a
localizer is needed that outputs a per-pair signal. This section builds
and tests the natural next attempt: ask the same frozen judge a pairwise
question directly.

\hypertarget{method}{%
\subsection{6.2 Method}\label{method}}

For each candidate tool pair (A, B), one judge call
(\texttt{llama3.1:8b}, frozen) asks whether a task intended for A could
plausibly select B instead, and vice versa. Ground truth is a
pre-registered, 24-pair fixture (4 CONFUSED / 20 NOT\_CONFUSED) built
entirely from already-collected behavioral trial data --- 6 fresh
servers from Section 5's EXP-1 scoring (14 pairs), plus \textbf{10}
pairs from the RW1/RW2 anchors (5 GitHub + 5 AWS IAM), including several
pairs the \emph{old} Levenshtein-based heuristic had already
false-positived (a deliberately adversarial inclusion). 3 trials per
pair (seed 42+trial\_idx), majority vote. The pre-committed bar for ``a
real positive method'': precision $\geq$ 0.50 \textbf{and} recall $\geq$ 0.50.

\hypertarget{binary-result}{%
\subsection{6.3 Binary result}\label{binary-result}}

Confusion matrix: TP=4, FP=20, FN=0, TN=0. \textbf{Precision = 0.167,
recall = 1.00} --- recall clears the bar; precision does not. All
\textbf{24 of 24} pairs were verdicted CONFUSABLE, including all
\textbf{ten} sampled pairs from the two 100\%-accuracy anchor servers
(GitHub, AWS IAM) and both pairs where the old heuristic already
false-positived. The judge catches every real behavioral confusion
(perfect recall) by flagging almost everything --- a direct yes/no
confusability question elicits a near-uniform ``yes'' from this judge on
this fixture, not a discriminating per-pair signal.

\hypertarget{graded-confidence-retry}{%
\subsection{6.4 Graded-confidence retry}\label{graded-confidence-retry}}

One further, pre-registered, time-boxed retry (ratified before running,
per the ``is this a framing artifact'' question) replaced the binary
question with a 0--10 confidence score, using the same aggregation and
ground truth, threshold $\geq$5.0 (scale midpoint, fixed before any result
was seen). Result: \textbf{every one of the 24 pairs'} mean score landed
in a \textbf{5.00--5.67 band} --- identical confusion matrix to the
binary run (precision 0.167, recall 1.00). This is a \emph{different}
degenerate failure mode from the binary run's uniform-yes: rather than
defaulting to one categorical answer, the judge anchors near the numeric
midpoint of its own scoring guide regardless of input --- genuinely
confused pairs and 100\%-resolved anchor pairs score in the same narrow
band.

\hypertarget{conclusion}{%
\subsection{6.5 Conclusion}\label{conclusion}}

The negative is robust to framing, not an artifact of question phrasing:
two structurally different question formats (categorical yes/no; graded
0--10 confidence) produce the same confusion matrix via two different
degeneracy mechanisms. \textbf{Pairwise judging, binary or graded, with
this frozen judge, is not a usable per-pair confusability signal on this
fixture.} Per pre-registration, this is a hard stop --- no third variant
was attempted. We note the scope of what this rules out: the experiment
validates the \emph{judging} step only, given a candidate pair; it does
not test an automatic candidate-\emph{generation} mechanism (which pairs
to even consider from a full catalog), and two of the four real
confusions in the ground truth cross mechanical name-prefix family
boundaries that a family-scoped candidate generator would not have
proposed in the first place.

\textbf{This is itself a boundary finding, not a second defeat stacked
on Section 4's negatives.} Section 4 established that the helps-regime
is real and characterizable at the catalog level (density, source
access, headroom). This section establishes a further, independent
boundary: even where the regime is real, it is not currently
\emph{findable} below the cost of running the full behavioral protocol,
using the one localization approach natural to try with an existing
frozen judge. That is new information the map did not have before this
experiment --- it narrows where a cheap pre-deployment check can
substitute for the protocol, which is a distinct question from whether
the regime itself exists.

\begin{center}\rule{0.5\linewidth}{0.5pt}\end{center}

\hypertarget{discussion}{%
\section{7. Discussion}\label{discussion}}

\hypertarget{synthesis}{%
\subsection{7.1 Synthesis}\label{synthesis}}

Read together, Sections 4--6 answer the paper's title question precisely
rather than broadly. Tool-description quality \emph{does} change agent
selection behavior --- but the evidence for this comes from separate
findings that compose, not one single joint condition: an oracle-effect
bounded to one tested catalog density (§4.2.1), a generation-safety
condition that requires documented source (§4.2.2), and a prevalence
measurement, on the general behavioral regime construct (§5.1), that a
pre-registered pilot sample of real public Python servers places at 0 of
9 tested (Section 5).
Separately, even where the regime does occur, it has not been shown
findable cheaply: a pairwise LLM-judge localizer fails to identify which
specific tool pairs are confusable on a new, unscored server (Section
6). Outside the helps-conditions, the same intervention that helps
direct selection in the tested fixtures can be neutral, harmful to
selection on at least one already-resolved family, or harmful to
retrieval on the one synthetic catalog tested --- three
independently-measured non-regimes, not one, and not (yet) shown to
generalize beyond the fixtures each was measured on.

\hypertarget{a-diagnostic-practitioners-can-use-today-the-two-condition-regime-test}{%
\subsection{7.2 A diagnostic practitioners can use today: the
Two-Condition Regime
Test}\label{a-diagnostic-practitioners-can-use-today-the-two-condition-regime-test}}

Section 5.1's regime definition is not only an analysis tool for this
paper --- the same two-step procedure is directly usable, unchanged, as
a pre-investment check on a given MCP server:

\begin{quote}
\textbf{The Two-Condition Regime Test}

\begin{enumerate}
\def\labelenumi{\arabic{enumi}.}
\tightlist
\item
  \textbf{Fail} --- Does the agent fail at least one contested task
  under the server's real, currently-shipped descriptions?
\item
  \textbf{Recover} --- If it fails, does a hand-written, ground-truth
  (oracle) description recover it?
\end{enumerate}

\begin{itemize}
\item
  (1) is NO $\rightarrow$ the agent already resolves the task from context;
  description tooling has nothing to fix here.
\item
  (1) is YES but (2) is also NO $\rightarrow$ the failure is not description-shaped
  (e.g., genuine functional overlap between tools, or a
  catalog-overwhelm failure mode); description tooling will not help
  either.
\item
  (1) and (2) are both YES $\rightarrow$ this server's family is inside the regime
  this paper's Section 4.2 findings speak to.
\end{itemize}
\end{quote}

This is exactly the check used to produce every OUT-OF-REGIME/IN-REGIME
verdict in Section 5's table --- the check \emph{this paper used}, not a
separately validated general instrument. It has been exercised on ten
Python MCP servers plus two anchors (Section 5); nothing about the
check's own reliability has been tested on a larger or non-Python
population, and naming it is meant to make it adoptable, not to claim it
has been validated beyond what this pilot supports. Applying it does not
require re-running this paper's full protocol at T18 scale --- steps (1)
and (2) are exactly the frozen-protocol headroom gate and oracle A/B
(§3), run at whatever scale the server in question actually has.

The practical implication for MCP server authors and tool-catalog
builders is not ``descriptions don't matter'' --- it is ``check whether
you are in the regime before investing in description tooling,'' where
``the regime'' is checked behaviorally via the test above, rather than
assumed from catalog size alone. A server whose agent already resolves
tasks from context gains nothing from more precise descriptions and
risks the account\_query-style harm pattern (Section 4.3.2) if it
applies them blanket. A server that \emph{is} in the regime should
generate descriptions from documented source with a target-grounded,
non-comparative prompt (Section 4.2.2's Q5/Q6 configuration), not from
the interface alone.

\hypertarget{what-would-change-this-picture}{%
\subsection{7.3 What would change this
picture}\label{what-would-change-this-picture}}

A larger or non-Python-verified EXP-1 sample (Section 5.5); a true
proprietary-frontier-model replication of Section 4.2.3; or a localizer
signal that does not route through this frozen judge's pairwise-question
failure mode (Section 6.5) --- for example, a signal derived from
name-embedding distance plus schema overlap rather than an LLM judgment
call.

\begin{center}\rule{0.5\linewidth}{0.5pt}\end{center}

\hypertarget{threats-to-validity-limitations}{%
\section{8. Threats to Validity /
Limitations}\label{threats-to-validity-limitations}}

Full list: \url{docs/paper/threats_to_validity.md}. Summarized here
by category; every item there traces to a specific finding already
reported in Sections 4--6, not a hedge added at write-up time.
\textbf{The note opening §8 is read first} --- it is the single most important threat
to a paper whose headline results are nulls, and it is stated
standalone, not as a clause inside another paragraph.

\hypertarget{the-false-negative-asymmetry-read-this-one-first}{%
\subsection{The false-negative asymmetry --- the epistemic bound on this
paper's null claims (read this one
first)}\label{the-false-negative-asymmetry-read-this-one-first}}

\textbf{Every null and boundary claim in this paper --- EXP-1's 0-of-9
headline, EXP-3's localizer-fails headline --- is bounded by one
asymmetry in what this pipeline's error-detection can catch, and this
section states that bound precisely.} Section 5.4 reports two
false positives from a seed-configuration bug, caught and reversed
before publication \emph{because} they were surprising enough to trigger
a recheck (a server unexpectedly showing in-regime behavior got a second
look). That same mechanism has a mirror-image blind spot: a bug that
instead silently \emph{suppressed} a real in-regime
signal --- turning a true positive into a false null --- would produce a
result indistinguishable from a correctly-measured null: ``this server
is not in-regime'' looks the same whether it's true or an artifact, and
nothing about a null result prompts the same ``that's surprising, let's
recheck'' response that caught the two false positives. We have not
found evidence of such a bug; we also have no mechanism that would
necessarily surface one, which is precisely the point --- this is the
bound this paper's nulls should be read against, not a suggestion that
they are unreliable. Any reader treating EXP-1 or EXP-3's null as more
secure than ``we looked and did not find a positive, using a pipeline
whose false-positive-catching capacity is demonstrated and whose
false-negative-catching capacity is not'' should recalibrate against
this paragraph specifically.

\hypertarget{sampling-generalizability}{%
\subsection{8.1 Sampling /
generalizability}\label{sampling-generalizability}}

N=10, Python-only, public-GitHub pilot. The strength claimed is
\emph{convergence} --- the EXP-1 null, the RW1/RW2 anchors, the P2-A
account\_query harm, and the doc-density-rarity pattern across tiers all
point the same direction independently --- not N=10 in isolation.
Non-Python mechanical extraction was dropped as a capability limit of
the method, not a data-availability inconvenience. Tier labels
(well\_documented/thin/near\_empty) are relative ranks within the
sampled pool, not absolute bands.

\hypertarget{agent-model-scope}{%
\subsection{8.2 Agent / model scope}\label{agent-model-scope}}

One default agent (gemma2:9b) across the regime map, Q-series, RW1/RW2,
and P2-A. The controlled capability ladder that would test
agent-capability sensitivity directly (EXP-2) was dropped from this
paper's scope --- ratified, justified by the regime already being
uncommon in the sampled population (a ladder over a rare regime has
limited external relevance). FRONTIER-T18 is one model, one data point,
not apples-to-apples with the gemma2:9b harness, and not a
proprietary-frontier-model result.

\hypertarget{measurement-judge-validity}{%
\subsection{8.3 Measurement / judge
validity}\label{measurement-judge-validity}}

Two seed-bug false positives were caught and reversed before reporting
(Section 5.4; asymmetry discussed prominently in §8's opening note, not here).
EXP-3's negative is robust across two framings, not a single-question
artifact, but validates the judging step only, not candidate-pair
generation. Trial counts deviate from the 5-per-arm default for
FRONTIER-T18 and EXP-3 (3 each), independently pre-registered and
justified.

\hypertarget{harm-asymmetric-risk}{%
\subsection{8.4 Harm / asymmetric risk}\label{harm-asymmetric-risk}}

Blanket description-fixing is not universally safe: P2-A's
account\_query family shows reproducible −20pp harm under two
independent generation strategies. Do-no-harm testing (Q6) covers only
the case where honest descriptions remain semantically distinct across
sibling tools; total description collapse to identical phrasing is
explicitly untested.

\hypertarget{reproducibility-artifact-note}{%
\subsection{8.5 Reproducibility-artifact
note}\label{reproducibility-artifact-note}}

The FRONTIER-T18 result (Section 4.2.3) required, during this paper's
preparation, recovering a raw result file that existed only as a
gitignored local artifact and independently re-deriving its numbers
before committing it as a hashed fixture (Section 9.2). This is now
resolved for the reported number; the harness code that produced it
is now merged to \texttt{main} as well (Section 9.2).

\begin{center}\rule{0.5\linewidth}{0.5pt}\end{center}

\hypertarget{reproducibility-artifact}{%
\section{9. Reproducibility Artifact}\label{reproducibility-artifact}}

\hypertarget{claim}{%
\subsection{9.1 Claim}\label{claim}}

Every figure in Sections 4--6 traces to a committed file on this
paper-writing branch, verified against a specific commit reachable from
\texttt{HEAD} (full trace: \url{docs/paper/evidence_table.md}). Most
fixture files carry SHA-256{[}:12{]} hashes recorded at pre-registration
time, before their experiment's results were known (protocol appendix:
\url{docs/research/frozen_protocol.md}). \textbf{The FRONTIER-T18
fixtures are the stated exception} --- see §9.2 immediately below: their
hashes were computed post-hoc, during this paper's preparation, not at
pre-registration time, because the original pre-registered artifact had
to be recovered from a gitignored local file rather than read from a
tracked one.

\hypertarget{frontier-t18-datacode-split-state-plainly-not-silently}{%
\subsection{9.2 FRONTIER-T18 data/code split (state plainly, not
silently)}\label{frontier-t18-datacode-split-state-plainly-not-silently}}

The result data underlying Section 4.2.3 is committed and independently
re-derivable:
\url{evals/fixtures/frontier_t18_step2_result.json} (per-task,
sha256{[}:12{]}=\texttt{3ca4a25dbd25}) and
\url{evals/fixtures/frontier_t18_step2_raw_calls.json} (all 240
raw per-call LLM responses, sha256{[}:12{]}=\texttt{93fb0d77262d}) ---
re-counting \texttt{SELECTED-CORRECT} directly from the raw calls
reproduces 71/120 (Arm A) and 120/120 (Arm B) exactly. The
\textbf{harness code} that produced these files
(\url{agentgauge/frontier.py},
\url{scripts/run_frontier_t18.py}) is now committed to \texttt{main}
alongside the data --- a from-scratch re-run of this experiment is
possible directly from this repository, with no further merge required.

\hypertarget{governance-1}{%
\subsection{9.3 Governance}\label{governance-1}}

All research-program branches route through draft, human-reviewed PRs;
any change touching the judge, scorer, rubric, or calibration constants
is escalated before merge (Section 3.1). No PR in this research program
has auto-merge; a human merges every one.

\begin{center}\rule{0.5\linewidth}{0.5pt}\end{center}

\hypertarget{conclusion-1}{%
\section{10. Conclusion}\label{conclusion-1}}

Tool-description quality changes agent tool-selection behavior in a
real, measurable, and mechanistically legible way --- but only inside a
narrow, checkable set of conditions: a catalog at least as dense as the
one point where an effect was tested and found (§4.2.1), and generation
with access to documented source. That effect does not collapse under a
substantial capability increase in the one additional agent tested
(Section 4.2.3) --- a finding \emph{inconsistent with} ``this is a
gemma2:9b-specific capability artifact,'' though one data point on a
different harness does not prove the effect is capability-independent in
general. Outside these conditions: the same behavioral regime construct
(§5.1) appears in 0 of 9 servers with a testable confusable family in a
pre-registered pilot sample of real, public, documented Python MCP
servers (Section 5); the same kind of intervention harms selection on at
least one
already-resolved family tested, while two other already-resolved
families in the same experiment show no harm (Section 4.3.2); it harms
retrieval across every retriever type tested on the one synthetic
catalog used for that test (Section 4.3.3); and the regime cannot
currently be located cheaply by asking an LLM judge pairwise whether two
tools are confusable (Section 6) --- itself a boundary finding, not a
second defeat (§6.5). One threat bounds all of these null and boundary
claims at once and is not a footnote: a bug that silently suppressed a
real in-regime signal would look identical to a correctly-measured null,
and this paper's demonstrated error-detection catches false positives,
not false negatives (§8, opening note). With that caveat stated plainly, the
field's assumption that description-quality improvement is a broadly
beneficial, low-risk lever reads, on this evidence, as over-general
rather than wrong: this paper's contribution is a first pass at the
boundary that makes ``broadly'' precise enough to check and re-test.

\begin{center}\rule{0.5\linewidth}{0.5pt}\end{center}

\hypertarget{appendix}{%
\section{Appendix}\label{appendix}}

\begin{itemize}
\item
  \textbf{A.1} Full cross-experiment regime table:
  \url{docs/research/exp4_regime_map.md} §Cross-Experiment Regime
  Table.
\item
  \textbf{A.2} Frozen judge/generator prompts (Guard-B target-grounded
  prompt; pairwise binary/graded confusability prompts):
  \url{agentgauge/fixer.py}
  (\url{_DESC_GENERATOR_GUARD_B_PROMPT}),
  \url{docs/research/exp3_pre_registration.md} §2, §7.
\item
  \textbf{A.3} Fixture hash manifest: generate per
  \url{docs/research/frozen_protocol.md} Appendix; the two
  FRONTIER-T18 hashes added during this paper's preparation are recorded in Section 9.2
  above and in \url{docs/paper/evidence_table.md} §1.3.
\item
  \textbf{A.4} Per-server EXP-1 table: Section 5.3 above (reproduced
  from \texttt{STATUS.md} EXP-1 section).
\item
  \textbf{A.5} Bibliography --- every citation below was independently
  fetched and checked against its claimed content during this paper's preparation (title,
  authors, and the specific figure/quote cited all confirmed against the
  primary source, not copied from a prior internal desk-research doc
  without re-verification). None were invented; where a candidate
  citation from prior internal research
  (\url{reports/frontier_phase1_research.md}, itself
  gitignored/uncommitted and self-flagged as ``arXiv IDs not
  independently re-resolved'') could not be verified to support its
  claimed content, it was dropped rather than included on trust --- see
  the exclusions list below.

  \textbf{Cited in this paper:}

  \begin{enumerate}
  \def\labelenumi{\arabic{enumi}.}
  \tightlist
  \item
    Anthropic Engineering, ``Writing tools for agents'' (publication
    date not independently confirmed during this paper's preparation --- content verified,
    date not).
    \url{https://www.anthropic.com/engineering/writing-tools-for-agents}
    --- verified: contains the quoted claim that small tool-description
    refinements yielded SOTA SWE-bench Verified results for Claude
    Sonnet 3.5. \nocite{anthropic2026tools}
  \item
    GitHub Engineering, ``Measuring what matters: how offline evaluation
    of GitHub MCP Server works'' (publication date not independently
    confirmed during this paper's preparation --- content verified, date not).
    \url{https://github.blog/ai-and-ml/generative-ai/measuring-what-matters-how-offline-evaluation-of-github-mcp-server-works/}
    --- verified: contains the quoted claim that small description edits
    ``shift results a lot.'' \textbf{Second-pass audit note:} this is a
    claim about the \emph{same} GitHub MCP server mirrored in RW1
    (§4.3.1), which found zero headroom for gemma2:9b on that server.
    The tension is flagged explicitly in §1.1 body text, not silently
    juxtaposed; not resolved (harness/model/metric details of GitHub's
    own eval were not verified during this paper's preparation). \nocite{github2026measuring}
  \item
    Hasan, M. M., Li, H., Rajbahadur, G. K., Adams, B., and Hassan, A.
    E., ``Model Context Protocol (MCP) Tool Descriptions Are Smelly!
    Towards Improving AI Agent Efficiency with Augmented MCP Tool
    Descriptions,'' arXiv:2602.14878, 2026. Verified via
    \url{arxiv.org/abs/2602.14878} --- 856 tools / 103 servers,
    97.1\% $\geq$1 ``smell,'' 56\% unclear purpose, +5.85pp median
    task-success gain from augmentation, 16.67\% regression rate,
    +67.46\% execution-step inflation, all confirmed against the
    abstract. \nocite{hasan2026smelly}
  \item
    Lu, X., Huang, H., Meng, R., Jin, Y., Zeng, W., and Shen, X.,
    ``Tools are under-documented: Simple Document Expansion Boosts Tool
    Retrieval'' (Tool-DE), arXiv:2510.22670, 2025 (arXiv ID year prefix
    ``25'' = 2025 --- corrected in second-pass audit; originally
    mislabeled 2026). Verified via \url{arxiv.org/abs/2510.22670} ---
    Tool-Embed/Tool-Rank, SOTA retrieval gains from LLM-driven document
    expansion, confirmed against the abstract. \nocite{lu2025toolde}
  \item
    Shi, Z., Wang, Y., Yan, L., Ren, P., Wang, S., Yin, D., and Ren, Z.,
    ``Retrieval Models Aren't Tool-Savvy: Benchmarking Tool Retrieval
    for Large Language Models'' (ToolRet), ACL 2025 Findings
    (\url{aclanthology.org/2025.findings-acl.1258}) /
    arXiv:2503.01763. Verified via web search cross-referencing the ACL
    Anthology and arXiv listings --- 7.6k tasks / 43k tools, best model
    (NV-embed-v1) 33.83 nDCG@10, both figures confirmed. \nocite{shi2025toolret}
  \item
    Yu, P., Liu, W., Yang, Y., Li, J., Zhang, Z., Feng, X., and Zhang,
    F., ``Benchmarking LLM Tool-Use in the Wild'' (WildToolBench),
    arXiv:2604.06185, 2026. Verified via
    \url{arxiv.org/abs/2604.06185v1} --- confirmed: ``no model
    achieves an accuracy of more than 15\%'' across 57 evaluated LLMs.
    \textbf{Not cited: a specific ``Grok-4 24.07\% wrong-name error''
    figure} from the prior internal desk-research doc --- this could not
    be confirmed from the abstract and full-text confirmation was not
    obtained; excluded rather than asserted on trust. \nocite{yu2026wildtool}
  \item
    Babu, R. S., and Iyer, L. G., ``ToolChoiceConfusion: Causal Minimal
    Tool Filtering for Reliable LLM Agents,'' arXiv:2606.06284, 2026.
    Verified via \url{arxiv.org/abs/2606.06284v1} --- confirmed:
    Causal Minimal Tool Filtering (CMTF), 102 tasks / 100 tools / four
    LLM backends, reliability-vs-tool-menu-size problem statement.
    \textbf{Not cited: the prior internal doc's characterization of this
    paper's evaluation as ``fully synthetic'' with an explicit
    disclaimer of ``no empirical measurement of how frequently this
    confusion regime occurs in practice''} --- this specific
    framing/quote could not be confirmed from the fetched abstract;
    excluded rather than asserted on trust. \nocite{babu2026toolchoice}
  \item
    GitHub issue, ``Duplicate tool names across MCP servers cause
    errors,'' \url{github.com/openai/openai-agents-python#464}
    (closed). Verified via \texttt{gh\ api} --- title and closed state
    confirmed. \nocite{openai2024issue464}
  \item
    GitHub issue, ``Feedback for dynamic tool selection,''
    \url{github.com/github/github-mcp-server#275} (closed). Verified
    via \texttt{gh\ api} --- title and closed state confirmed;
    referenced only in this appendix as supporting context, not cited in
    body prose in this draft. \nocite{github2024issue275}
  \end{enumerate}

  \textbf{Checked and excluded (candidate citations from prior internal
  desk research that did NOT verify --- listed so the exclusion is
  visible, not silent):}

  \begin{itemize}
  \tightlist
  \item
    ``arXiv:2603.20313, retrieval quality fundamentally bounded by the
    informativeness of tool descriptions'' --- fetched; the paper at
    this ID (``Semantic Tool Discovery for Large Language Models: A
    Vector-Based Approach to MCP Tool Selection'') does not support this
    claim per its abstract. Not used anywhere in this paper.
  \item
    Any secondary-source selection-accuracy-vs-tool-count percentages
    (e.g., ``\textgreater90\% at 5-7 tools → \textasciitilde13\% at 100+
    tools'') from \url{docs/research/phase1-buyer-and-landscape.md}
    --- that document's own ``Confidence \& gaps'' section already flags
    these as ``secondary-source paraphrase, not verified against a
    primary benchmark.'' Not used in this paper for the same reason its
    original author already gave.
  \end{itemize}
\item
  \textbf{A.6} EXP-1 sampling-frame revision history (moved here from
  §5.2 body --- the finding that survives
  in §5.2 is unaffected; this is provenance, not evidence). The
  pre-registered frame was rebuilt five times over the course of the
  experiment, each rebuild escalated and ratified before proceeding,
  never silently:

  \begin{itemize}
  \tightlist
  \item
    \textbf{v1} --- star-stratified GitHub-topic pool.
  \item
    \textbf{v2} --- doc-density-stratified, N=23.
  \item
    \textbf{v3--v4} --- three, then a fourth, confirmed
    Python-AST-extractor bug found and fixed while preparing the trial
    batch.
  \item
    \textbf{v5 (final, ratified)} --- a systematic audit found that the
    generic regex fallback used for every non-Python server pulls in
    systemic noise (template literals, parameter names, category labels,
    unrelated example data) --- a capability limit of blind
    text-proximity matching, not a fixable parsing bug. All 11
    non-Python servers were dropped, yielding the final N=10,
    Python-only frame reported in §5.2. Full commit-level trail:
    \texttt{STATUS.md} EXP-1 section, ``Frame history'' paragraph;
    ratification commit \texttt{538affe};
    \url{docs/paper/evidence_table.md} §2.
  \end{itemize}
\item
  \textbf{A.7} Full Q3→Q6 source-aware generation progression (moved
  here from §4.2.2 body --- both prose
  findings in §4.2.2 are unaffected; this is intermediate detail, not a
  third finding). The two-row table in §4.2.2 shows only the endpoints
  (interface-only failure; Q5/Q6 safe-and-recovering). The full
  four-stage progression, showing \emph{how} Q5/Q6 was reached and the
  two distinct fabrication failure modes it closes:

  \renewcommand{\arraystretch}{1.15}
  \begin{longtable}{@{}p{0.30\linewidth}p{0.19\linewidth}p{0.11\linewidth}p{0.26\linewidth}@{}}
  \toprule
  \textbf{Condition} & \textbf{Recovery (6 structural contested tasks)} & \textbf{p} &
  \textbf{No-fabrication} \\
  \midrule
  \endhead
  Q3 F-DOC (whole-file source + docstrings) & 83.3\% (5/6) & 0.0625
  (marginal) & PASS \\
  Q3 F-BODY (whole-file, docstrings stripped) & 83.3\% (invalidated) &
  --- & \textbf{FAIL} --- cross-tool source misattribution \\
  Q4-DOC-scoped (per-tool source + docstrings) & 100\% (6/6) & 0.0313 &
  \textbf{FAIL} --- docstring-body vocabulary gap \\
  Q4-BODY-scoped & 100\% (6/6) & 0.0313 & PASS \\
  Q5-guarded (per-tool source + docstrings + target-grounded prompt) &
  100\% (6/6) & 0.0313 & PASS \\
  Q6-guarded (23-tool blanket application) & 100\% (6/6) & 0.0312 &
  PASS, \textbf{0/11 regressions} on already-passing tools \\
  \bottomrule
  \end{longtable}
  \renewcommand{\arraystretch}{1}

  Reading order: Q3 (whole-file) shows docstrings are safe there but
  only marginally recovering, and body-only is unsafe (cross-tool
  misattribution). Q4 (scoped) fixes the cross-tool issue but
  \emph{inverts} the safety finding --- scoped docstrings now fabricate
  via a docstring-body vocabulary gap, while scoped body-only is safe
  but discards the docstring signal. Q5 (guarded prompt) resolves the
  inversion: scoped + docstrings + a target-grounded, non-comparative
  prompt is safe AND 100\% recovering. Q6 confirms this holds blanket
  across a larger, 23-tool catalog with zero regressions. Full
  narrative: \url{docs/research/exp4_regime_map.md} §Regime 2;
  \texttt{STATUS.md} Q3/Q4/Q5/Q6 sections;
  \url{docs/paper/evidence_table.md} §1.2.
\end{itemize}


\bibliographystyle{plainnat}
\bibliography{references}

\end{document}
